The algebraic structures introduced in Sections~\ref{sec:casimir} and~\ref{sec:YB}
provide the general framework for the $r$-matrix formalism.
We now show how these structures are realized concretely in the Toda chain, which gives a fundamental example of an integrable system associated with a simple Lie algebra \cite{toda1967vibration,Olshanetsky1979,KOSTANT1979195,perelomov1990integrable}.
\subsection{The model}\label{subsec:The model}
The open Toda chains are integrable systems associated with Lie algebras.
Let $\mathcal{G}$ be a finite-dimensional simple Lie algebra of rank $r$, equipped with an invariant scalar product denoted by $( \cdot, \cdot )$. We choose a Cartan subalgebra with an orthonormal basis $\{H_i\}$. We have the Cartan decomposition of $\mathcal{G}$:
\begin{equation}
    \mathcal{G} = \mathcal{N}_{-} \oplus \mathcal{H} \oplus \mathcal{N}_{+} \quad \text{with} \quad \mathcal{N}_{\pm} = \bigoplus_{\beta \text{ positive}} \mathcal{G}_{\pm \beta},
\end{equation}
where $\mathcal{G}_{\pm \beta}$ are one-dimensional, generated by the root vectors $E_{\pm \beta}$ for any positive root $\beta$.
We will need the commutation relations:
\begin{align}
    [H_i, H_j] &= 0, \\
    [H, E_{\pm \alpha}] &= \pm \alpha(H) \, E_{\pm \alpha}, \\
    [E_{\alpha}, E_{-\alpha}] &= ( E_{\alpha}, E_{-\alpha} ) H_\alpha,
\end{align}
where $H_{\alpha} = \sum_i \alpha(H_i) \, H_i$. We will often use the constants $n_\alpha$ such that 
\begin{equation}
    n_\alpha^2 ( E_{\alpha},E_{-\alpha} ) = 1.
\end{equation}
Toda chains are associated with any simple Lie algebra $\mathcal{G}$ as follows. Consider two elements in the Cartan subalgebra:
\begin{equation}
    q = \sum_{i = 1}^r q_i H_i, \quad p = \sum_{i = 1}^r p_i H_i \quad r = \rank \mathcal{G}.
\end{equation}
The coefficients $(q_i, p_i)$ are the coordinates of a $2r$-dimensional phase space with Poisson brackets 
\begin{equation}\label{eq:poison_bis}
    \{q_i, p_j\} = \frac{1}{2} \delta_{ij},
\end{equation}
where the factor $1/2$ is introduced to simplify later formulae.
The Hamiltonian of the open Toda chain is by definition:
\begin{equation}
    H = ( p, p ) + 2 \sum_{\alpha \text{ simple}} \mathrm{e}^{2\alpha(q)}.
\end{equation}
The equations of motion are then:
\begin{align}
    \dv{q}{t} &= \{ q , H \} = p \\
    \dv{p}{t} &= \{ p , H \} = - 2 \sum_{\alpha \text{ simple}} H_\alpha \mathrm{e}^{2\alpha(q)}
\end{align}
or, combining the two,
\begin{equation}
    \dv[2]{q}{t} = - 2 \sum_{\alpha \text{ simple}} H_\alpha \mathrm{e}^{2\alpha(q)}.
\end{equation}

In the usual $\mathfrak{sl}(r+1)$ case, the Lie algebra of traceless $(r+1) \times (r+1)$ matrices, the Cartan algebra can be taken as traceless diagonal matrices.
The simple root vectors $E_{\alpha_i}$ are the canonical matrices $E_{i, i+1}$.
In this case it is not very convenient to use an orthonormal basis of the Cartan algebra under the scalar product $ ( X,Y )  = \Tr(XY)$.
Instead we write an element in the Cartan subalgebra $q = \sum_{i=1}^{r+1} q_i E_{ii}$ with the constraint $\sum q_i = 0$. 
In this description the simple roots are 
\begin{equation}
    \alpha_i(q) = q_i - q_{i+1}.
\end{equation}
The dynamical variables of the open Toda chain are the two elements $q = \sum_{i = 1}^{r+1} q_i E_{ii}$ and $p~=~ \sum_{i=1}^{r+1}p_iE_{ii}$ with the constraints $\sum_i q_i = 0$ and $\sum_i p_i = 0$. These constraints are not compatible with the canonical Poisson bracket, eq. (\ref{eq:Foundamental_Poisson}), so we use instead the Dirac bracket:
\begin{equation}
    \{ q_i, p_j \} = \frac{1}{2} \left( \delta_{ij} - \frac{1}{r+1} \right), \quad \{ p_i, p_j \} = \{ q_i, q_j \} = 0.
\end{equation}
The Hamiltonian 
\begin{equation}
    H = \sum_{i=1}^{r+1} p_i^2 + 2 \sum_{i=1}^r \mathrm{e}^{2(q_i - q_{i+1})}
\end{equation}
generates the equations of motion:
\begin{align}
    \dot{q}_i &= p_i \\
    \dot{p}_1 &= - 2 \mathrm{e}^{2(q_1 - q_2)} \\
    \dot{p}_i &= - 2 \mathrm{e}^{2(q_{i} - q_{i+1})} + 2 \mathrm{e}^{2(q_{i-1} - q_{i})} \\
    \dot{p}_{r+1} &= 2 \mathrm{e}^{2(q_r - q_{r+1})}.
\end{align}
The particular form of the equations for $\dot{p}_1$ and $\dot{p}_{r+1}$ explains the terminology \textit{open} Toda chain.

The equations of motion of the open Toda chain can be written in Lax form and the system is integrable.
Note that the exponential dependence in the Lax matrix is chosen so that the commutator $[M,L]$ produces terms proportional to $e^{2\alpha(q)}$, matching the potential term in the Hamiltonian.

\begin{proposition}
    The equations of motion of the open Toda chain admit a Lax pair representation $\dot{L} = {[M,L]}$ with 
    \begin{equation} \label{eq:lax_matrix_toda_chain}
        L = p + \sum_{\alpha \text{ simple}} n_\alpha \mathrm{e}^{\alpha(q)} (E_\alpha + E_{-\alpha}), \quad M = - \sum_{\alpha \text{ simple}} n_\alpha \mathrm{e}^{\alpha(q)} (E_\alpha - E_{-\alpha}).
    \end{equation}
\end{proposition}
\begin{proof}
   We compute the time derivative of $L$:
   \begin{equation}
       \dv{L}{t} = \dv{p}{t} + \sum_\alpha n_\alpha \alpha \left( \dv{q}{t} \right) \mathrm{e}^{\alpha(q)} (E_\alpha + E_{-\alpha}).
   \end{equation}
   Then we compute the commutator
   \begin{equation}
       [M,L] = \sum_\alpha n_\alpha \alpha(p) \mathrm{e}^{\alpha(q)} (E_\alpha + E_{- \alpha}) - \sum_\alpha \sum_\beta n_\alpha n_\beta \mathrm{e}^{(\alpha + \beta)(q)} [E_\alpha - E_{-\alpha}, E_\beta + E_{-\beta}].
   \end{equation}
   Since $[E_{\pm \alpha}, E_{\pm \beta}]$ is antisymmetric in the exchange of $\alpha$ and $\beta$, the second term reduces to 
   \begin{equation}
    - \sum_\alpha \sum_\beta n_\alpha n_\beta \mathrm{e}^{(\alpha + \beta)(q)} ([E_\alpha, E_{-\beta}] - [E_{-\alpha},E_{\beta}]).
   \end{equation}
   Since $\alpha$ and $\beta$ are simple roots, we have $[E_\alpha, E_{-\beta}] = \delta_{\alpha\beta}( E_{\alpha}, E_{-\alpha} ) H_\alpha$
   and therefore
   \begin{equation}
       [M,L] = -2\sum_\alpha \mathrm{e}^{2\alpha(q)} H_\alpha + \sum_{\alpha} n_\alpha \alpha(p) \mathrm{e}^{\alpha(q)} (E_\alpha + E_{-\alpha})
   \end{equation}
   which reproduces exactly the equations of motion, hence proving the Lax representation.
\end{proof}


\subsection[The \textit{r}-matrix of the Toda models]{The $\boldsymbol{r}$-matrix of the Toda models}\label{sec:toda_rmatrix}
The $r$-matrix of the Toda chain provides a canonical example of a constant $r$-matrix associated with a simple Lie algebra, and is known to satisfy the modified classical Yang–Baxter equation. In this section, we derive its explicit form and verify the associated Poisson structure.

\begin{theorem}
    Let $L$ be the Lax matrix $(\ref{eq:lax_matrix_toda_chain})$, with $p$ and $q$ satisfying the Poisson bracket relation $(\ref{eq:poison_bis})$.
    Then there exists an antisymmetric $r$-matrix 
    $r_{12} = - r_{21} \in \mathcal{G}^{\otimes 2}$,
    independent of the phase space variables $p$ and $q$, such that
    \begin{equation}
        \{ L_1 , L_2 \} = [ r_{12},\, L_1 + L_2 ].
    \end{equation}
    It is given by
    \begin{equation} \label{eq:r_12_definition}
        r_{12} = \frac12 \sum_{\alpha\ \text{positive}}
        \frac{E_{\alpha} \otimes E_{-\alpha} \;-\;
              E_{-\alpha} \otimes E_{\alpha}}
             {( E_{\alpha}, E_{-\alpha} )}.
    \end{equation}
    This formula implies that the conserved Hamiltonians are in involution.
\end{theorem}


The proof relies on the following auxiliary results.

\begin{lemma} \label{lemma:r_H}
    Let $r_{12}$ be defined by \eqref{eq:r_12_definition}.  
    Then
    \begin{equation}
        \left[r_{12},\,H\otimes1+1 \otimes H\right]=0
        \qquad \forall\, H \in \mathcal H.
    \end{equation}
\end{lemma}
\begin{proof}
    We directly compute the commutator:
    \begin{equation}
        [E_\alpha \otimes E_{-\alpha}, H_i \otimes 1 + 1 \otimes H_i] = [E_\alpha, H_i] \otimes E_{-\alpha} + E_{\alpha} \otimes [E_{-\alpha}, H_i].
    \end{equation}
    Using 
    \begin{equation}
        [H_i,E_\alpha] = \alpha(H_i) E_\alpha
        \qquad
        [H_i,E_{-\alpha}] = - \alpha(H_i) E_{-\alpha}
    \end{equation}
    we obtain
    \begin{equation}
        -\alpha(H_i) \, E_{\alpha} \otimes E_{-\alpha} + \alpha(H_i) \, E_{\alpha} \otimes E_{-\alpha} = 0.
    \end{equation}
    Similarly $ [E_{-\alpha}\otimes E_\alpha,\,H_i\otimes1+1 \otimes H_i] = 0 $. \\ 
    Thus every term in the sum commutes with $H_i\otimes1+1 \otimes H_i$.
\end{proof}

\begin{lemma} \label{lemma:r_E}
    Let $r_{12} \in \mathcal{G} \otimes \mathcal{G}$ be defined by equation~\eqref{eq:r_12_definition}. 
    Then, for every simple root $\alpha$, one has
    \begin{equation}
        [r_{12},\, E_{\pm \alpha} \otimes 1 + 1 \otimes E_{\pm \alpha}]
        =
        \frac{1}{2}
        \bigl(
        H_{\alpha} \otimes E_{\pm \alpha}
        -
        E_{\pm \alpha} \otimes H_{\alpha}
        \bigr).
    \end{equation}
\end{lemma}
\begin{proof}
    We compute $ [r_{12}, E_\alpha\otimes1 + 1 \otimes E_{\alpha}]$.
    For simplicity, we suppress normalization factors depending on the roots. A fully explicit computation requires keeping track of the coefficients $(E_\beta, E_{-\beta})$, but does not affect the structure of the result.
    \begin{align}
        &= \frac12 \sum_{\beta \text{ positive}}
        [E_\beta\otimes E_{-\beta}-E_{-\beta}\otimes E_\beta,\,
        E_\alpha \otimes 1 + 1 \otimes E_\alpha ] \\
        &= \frac{1}{2} \sum_{\beta \text{ positive}}
        [E_\beta,E_\alpha]\otimes E_{-\beta}
        + E_\beta \otimes [E_{-\beta},E_\alpha]
        - [E_{-\beta},E_\alpha]\otimes E_\beta
        - E_{-\beta}\otimes [E_\beta,E_\alpha].
    \end{align}
    We divide the sum into two contributions: $\alpha = \beta$ and $\alpha \neq \beta$.
    Only the contribution with $\alpha=\beta$ produces Cartan elements, since $[E_\alpha,E_{-\alpha}] = H_\alpha$. We obtain:
    \begin{equation}
        \frac12
        \left(
        H_\alpha \otimes E_\alpha - E_\alpha \otimes H_\alpha
        \right).
    \end{equation}
    For every $\beta\neq\alpha$, the commutator $[r_{12}, E_{\pm\alpha} \otimes 1 + 1 \otimes E_{\pm\alpha}]$ is a linear combination of tensors of the form
    \begin{equation}
        E_{\beta\pm\alpha}\otimes E_{-\beta}, \qquad
        E_{\pm\alpha-\beta}\otimes E_{\beta}
    \end{equation}
    whenever these root spaces are nonzero. \\
    Consider the first type of terms: $E_{\beta + \alpha} \otimes E_{-\beta}$. We have
    \begin{align}
        [E_\beta \otimes E_{-\beta},\, E_\alpha \otimes 1] &= c_1 \, E_{\beta + \alpha} \otimes E_{-\beta} \\
        [E_{\alpha + \beta} \otimes E_{-\beta - \alpha},\, 1 \otimes E_{\alpha}] &= c_2 \, E_{\beta + \alpha} \otimes E_{-\beta}.
    \end{align}
    Hence,
    \begin{equation}
        [E_\beta \otimes E_{-\beta},\, E_\alpha \otimes 1] 
        + [E_{\alpha + \beta} \otimes E_{-\beta - \alpha},\, 1 \otimes E_{\alpha}]
        = (c_1 + c_2)\, E_{\beta + \alpha} \otimes E_{-\beta}.
    \end{equation}
   Using the invariance of the scalar product, i.e.
    \begin{equation}
        (X,[Y,Z]) = ([X,Y],Z),
    \end{equation}
    we compute
    \begin{align}
        c_1 &= (E_{-\beta-\alpha},[E_\beta,E_\alpha]), \\
        c_2 &= (E_\beta,[E_{-\beta-\alpha},E_\alpha])
            = -([E_\beta,E_\alpha],E_{-\beta-\alpha})
            = -c_1.
    \end{align}
    Therefore, the contributions of this type cancel, and their sum vanishes.
    Consider now the second type: $ E_{\alpha - \beta} \otimes E_\beta $.\\
    Since $\alpha$ is a simple root, whenever $\beta - \alpha$ is a root we may set $\gamma = \beta - \alpha$, so that $\beta = \alpha + \gamma$.
    Then
    \begin{equation}
        E_{\alpha - \beta} \otimes E_\beta = E_{-\gamma} \otimes E_{\alpha + \gamma}.
    \end{equation}
    This reduces the expression to the same form treated in the previous case, so the argument proceeds identically.

    Therefore, the only remaining terms are those with $\alpha = \beta$, which proves the claim:
    \begin{equation}
        [r_{12},\, E_{\pm \alpha} \otimes 1 + 1 \otimes E_{\pm \alpha}]
        = \frac{1}{2}\bigl(H_{\alpha} \otimes E_{\pm \alpha} - E_{\pm \alpha} \otimes H_{\alpha}\bigr).
    \end{equation}
\end{proof}
We now return to the proof of the theorem.
\begin{proof}[Proof of the theorem]
    We first compute the Poisson bracket of the Lax matrix:
    \begin{equation}
        \{L_1,L_2\}
        =
        \frac12 \sum_{\alpha\ \text{simple}}
        n_\alpha \mathrm{e}^{\alpha(q)}
        \bigl[
        H_\alpha \otimes (E_\alpha + E_{-\alpha})
        -
        (E_\alpha + E_{-\alpha}) \otimes H_\alpha
        \bigr].
    \end{equation}

    We now show that the tensor $r_{12}$ defined in \eqref{eq:r_12_definition} reproduces the same expression through the commutator $[r_{12},L_1+L_2]$.

    Since
    \begin{equation}
        L
        =
        p
        +
        \sum_{\alpha\ \text{simple}}
        n_\alpha \mathrm{e}^{\alpha(q)} (E_\alpha+E_{-\alpha}),
    \end{equation}
    we have
    \begin{equation}
        L_1+L_2
        =
        p \otimes 1 + 1 \otimes p
        +
        \sum_{\alpha\ \text{simple}}
        n_\alpha \mathrm{e}^{\alpha(q)}
        \bigl(
        (E_\alpha+E_{-\alpha}) \otimes 1
        +
        1 \otimes (E_\alpha+E_{-\alpha})
        \bigr).
    \end{equation}
    Therefore
    \begin{align}
        [r_{12},L_1+L_2]
        &=
        [r_{12},p \otimes 1 + 1 \otimes p] \notag \\
        &\quad +
        \sum_{\alpha\ \text{simple}}
        n_\alpha \mathrm{e}^{\alpha(q)}
        [r_{12},(E_\alpha+E_{-\alpha}) \otimes 1 + 1 \otimes (E_\alpha+E_{-\alpha})].
    \end{align}
    Since $p=\sum_i p_i H_i$, Lemma~\ref{lemma:r_H} implies
    \begin{equation}
        [r_{12},p \otimes 1 + 1 \otimes p]
        =
        \sum_i p_i [r_{12},H_i \otimes 1 + 1 \otimes H_i]
        =
        0.
    \end{equation}
    For the second term, by linearity of the commutator and by Lemma~\ref{lemma:r_E}, valid for every simple root $\alpha$, we obtain
    \begin{align}
        & [r_{12},(E_\alpha+E_{-\alpha}) \otimes 1 + 1 \otimes (E_\alpha+E_{-\alpha})] \notag \\
        &= [r_{12},E_\alpha \otimes 1 + 1 \otimes E_\alpha]
        +
        [r_{12},E_{-\alpha} \otimes 1 + 1 \otimes E_{-\alpha}] \notag \\
        &= \frac12
        \bigl(
        H_\alpha \otimes E_\alpha - E_\alpha \otimes H_\alpha
        \bigr)
        +
        \frac12
        \bigl(
        H_\alpha \otimes E_{-\alpha} - E_{-\alpha} \otimes H_\alpha
        \bigr) \notag \\
        &= \frac12
        \bigl[
        H_\alpha \otimes (E_\alpha+E_{-\alpha})
        -
        (E_\alpha+E_{-\alpha}) \otimes H_\alpha
        \bigr].
    \end{align}

    Substituting these identities into the expression of $[r_{12},L_1+L_2]$, we find
    \begin{equation}
        [r_{12},L_1+L_2]
        =
        \frac12 \sum_{\alpha\ \text{simple}}
        n_\alpha \mathrm{e}^{\alpha(q)}
        \bigl[
        H_\alpha \otimes (E_\alpha + E_{-\alpha})
        -
        (E_\alpha + E_{-\alpha}) \otimes H_\alpha
        \bigr].
    \end{equation}

    This coincides exactly with the expression of $\{L_1,L_2\}$. Hence
    \begin{equation}
        \{L_1,L_2\} = [r_{12},L_1+L_2].
    \end{equation}
    This proves the theorem.
\end{proof}


An alternative and more structural formulation of the $r$-matrix formalism can be obtained by introducing a linear operator $R:\mathcal G \to \mathcal G$ associated with $r_{12}$.
This point of view, which is closely related to the classical Yang--Baxter equation, is presented in Appendix~\ref{app:Rmatrix}.